\section{Stretched-exponential sums of ball kernels}
\label{sec:stretched_ball_kernel}

Let $I$ be finite, let $Y$ be an extended pseudometric space with
distance $\delta$, and let $i\mapsto x_i$ be arbitrary anchor
positions.  Use the ball kernels $K_r$ of
Definition~\ref{def:ball_incidence_kernel} and the moments $M_p$
of~(\ref{eq:stretched_moment}).
The following estimates give the geometric kernel bound in
\cite{OpenAI2026AreaLaw}, \texttt{09-amplification.tex},
lines 141--159.
% Source revision: adc7f1241b42e322a6451854ab7e4b4c146bf78a.
% Checked module: QICLean/Analysis/StretchedBallKernel.lean.
% This geometric package does not identify a physical oscillation channel
% or establish the full amplification theorem.

\begin{definition}[Stretched-exponential ball kernel]
\label{def:stretched_ball_kernel}
\lean{Metric.stretchedBallKernel}
\leanok
\uses{def:ball_incidence_kernel}
For real $C,c,\alpha$ and $y,z\in Y$, define
\begin{align}
 A_{C,c,\alpha}(y,z)
 &=\sum_{n=0}^{\infty}C e^{-cn^\alpha}K_n(y,z).
 \label{eq:stretched_ball_kernel}
\end{align}
Every radius $n$ is regarded as a finite nonnegative real radius.
The sum is the real infinite sum, with value zero if the series
does not converge.  Theorem~\ref{thm:stretched_ball_summable}
establishes convergence under the counting hypotheses used below.
\end{definition}

\begin{lemma}[Nonnegativity of the shell kernel]
\label{lem:stretched_ball_nonneg}
\lean{Metric.stretchedBallKernel_nonneg}
\leanok
\uses{def:stretched_ball_kernel}
For every $C\geq0$, $c,\alpha\in\mathbb R$, and $y,z\in Y$,
one has $A_{C,c,\alpha}(y,z)\geq0$.
\end{lemma}
\begin{proof}
\leanok
\uses{def:stretched_ball_kernel, lem:ball_incidence_nonneg}
Each summand is nonnegative by
Lemma~\ref{lem:ball_incidence_nonneg} and positivity of the
exponential.  Its infinite sum is nonnegative as well.
\end{proof}

\begin{lemma}[No entries across infinite distance]
\label{lem:stretched_ball_zero}
\lean{Metric.stretchedBallKernel_eq_zero_of_edist_eq_top}
\leanok
\uses{def:stretched_ball_kernel}
For every $C,c,\alpha\in\mathbb R$ and $y,z\in Y$, if
$\delta(y,z)=\infty$, then $A_{C,c,\alpha}(y,z)=0$.
\end{lemma}
\begin{proof}
\leanok
\uses{def:stretched_ball_kernel, lem:ball_incidence_zero}
Lemma~\ref{lem:ball_incidence_zero} gives $K_n(y,z)=0$ at every
natural radius $n$.  The series is identically zero.
\end{proof}

\begin{theorem}[Convergence from quadratic label counts]
\label{thm:stretched_ball_summable}
\lean{Metric.summable_stretchedBallKernel}
\leanok
\uses{def:ball_incidence_kernel}
Let $C\geq0$, $c>0$, $\alpha>0$, $B\in\mathbb R$, and fix
$y,z\in Y$.  Suppose that, for every $n\in\mathbb N$,
\begin{align}
 \bigl|\{i\in I:\delta(y,x_i)\leq n\}\bigr|
 &\leq B(1+n)^2.
 \label{eq:stretched_entry_labels}
\end{align}
Then the series
$\sum_{n=0}^{\infty}C e^{-cn^\alpha}K_n(y,z)$ is summable.
No finiteness of $Y$, site-count bound, or separate sign assumption
on $B$ is required; the $n=0$ label bound already implies $B\geq0$.
\end{theorem}
\begin{proof}
\leanok
\uses{lem:ball_incidence_nonneg, lem:ball_incidence_label_bound,
thm:stretched_moment_summable}
The summands are nonnegative by
Lemma~\ref{lem:ball_incidence_nonneg}.
Lemma~\ref{lem:ball_incidence_label_bound} and
(\ref{eq:stretched_entry_labels}) give, for every $n$,
\begin{align}
 \bigl|C e^{-cn^\alpha}K_n(y,z)\bigr|
 &\leq CB(1+n)^2e^{-cn^\alpha}.
 \label{eq:stretched_entry_majorant}
\end{align}
The right side is summable by
Theorem~\ref{thm:stretched_moment_summable} with $p=2$,
followed by multiplication by $CB$.
Comparison of norms proves convergence.  In particular, this is
absolute convergence of the actual shell series.
\end{proof}

\begin{theorem}[Uniform weighted row bound]
\label{thm:stretched_ball_weighted_row}
\lean{Metric.sum_stretchedBallKernel_mul_exp_le}
\leanok
\uses{def:stretched_ball_kernel}
Suppose $Y$ is finite.  Let $C,a,B,V\geq0$ and $\alpha>0$,
and let $c\in\mathbb R$ satisfy $\eta=c-a2^\alpha>0$.
Assume the following bounds for every $n\in\mathbb N$,
$y\in Y$, and $i\in I$:
\begin{align}
 \begin{aligned}
 \bigl|\{i\in I:\delta(y,x_i)\leq n\}\bigr|&\leq B(1+n)^2,\\
 \bigl|\{z\in Y:\delta(z,x_i)\leq n\}\bigr|&\leq V(1+n)^2.
 \end{aligned}
 \label{eq:stretched_ball_counts}
\end{align}
Then, for every $y\in Y$,
\begin{align}
 \begin{aligned}
 \sum_{z\in Y}A_{C,c,\alpha}(y,z)
 e^{a(\delta(y,z)_{\mathbb R})^\alpha}
 &\leq CBV\sum_{n=0}^{\infty}(1+n)^4
 e^{-(c-a2^\alpha)n^\alpha}\\
 &=CBV M_4(c-a2^\alpha,\alpha).
 \end{aligned}
 \label{eq:stretched_ball_weighted_row}
\end{align}
The remaining decay rate is exactly $c-a2^\alpha$.
The fourth moment is finite, and no assumption $\alpha\leq1$
is needed.
\end{theorem}
\begin{proof}
\leanok
\uses{def:stretched_ball_kernel, thm:stretched_ball_summable,
thm:ball_incidence_weighted_row, thm:stretched_moment_summable}
Since $a2^\alpha\geq0$, the hypothesis $\eta>0$ implies $c>0$.
Theorem~\ref{thm:stretched_ball_summable} proves convergence for
each entry.  Multiplying each entry's summands by the fixed
nonnegative weight $e^{a(\delta(y,z)_{\mathbb R})^\alpha}$
preserves convergence.  Since $Y$ is finite, the resulting finite
sum of convergent series is convergent, and its finite site sum
may be exchanged with the infinite radius sum.

At radius $n$, apply
Theorem~\ref{thm:ball_incidence_weighted_row} with the label and
site bounds $B(1+n)^2$ and $V(1+n)^2$ from
(\ref{eq:stretched_ball_counts}).  Multiplication by
$C e^{-cn^\alpha}\geq0$ gives
\begin{align}
 \begin{aligned}
 \sum_{z\in Y}C e^{-cn^\alpha}K_n(y,z)
 e^{a(\delta(y,z)_{\mathbb R})^\alpha}
 &\leq CBV(1+n)^4e^{-cn^\alpha}e^{a(2n)^\alpha}\\
 &=CBV(1+n)^4e^{-(c-a2^\alpha)n^\alpha}.
 \end{aligned}
 \label{eq:stretched_ball_shell_bound}
\end{align}
The equality uses $(2n)^\alpha=2^\alpha n^\alpha$ for
nonnegative bases, including $n=0$.
The majorant is summable by
Theorem~\ref{thm:stretched_moment_summable} with $p=4$ and
rate $\eta>0$.  Sum~(\ref{eq:stretched_ball_shell_bound}) over
$n$ using the established convergence, then factor out $CBV$.
This gives~(\ref{eq:stretched_ball_weighted_row}).
\end{proof}

\begin{theorem}[Explicit half-rate choice of the spatial weight]
\label{thm:stretched_ball_half_rate}
\lean{Metric.sum_stretchedBallKernel_mul_exp_half_rate_le}
\leanok
\uses{def:stretched_ball_kernel}
Suppose $Y$ is finite.  Let $C,B,V\geq0$, $c>0$, and
$\alpha>0$, and assume~(\ref{eq:stretched_ball_counts}) for all
$n\in\mathbb N$, $y\in Y$, and $i\in I$.
Set $a=c/(2\cdot2^\alpha)$.  Then $a>0$, and for every $y\in Y$,
\begin{align}
 \begin{aligned}
 \sum_{z\in Y}A_{C,c,\alpha}(y,z)
 e^{\frac{c}{2\cdot2^\alpha}(\delta(y,z)_{\mathbb R})^\alpha}
 &\leq CBV\sum_{n=0}^{\infty}(1+n)^4e^{-(c/2)n^\alpha}\\
 &=CBV M_4(c/2,\alpha).
 \end{aligned}
 \label{eq:stretched_ball_half_rate}
\end{align}
\end{theorem}
\begin{proof}
\leanok
\uses{thm:stretched_ball_weighted_row}
Since $2^\alpha>0$ and $c>0$, the chosen $a$ is positive.
Direct cancellation gives
$c-\bigl(c/(2\cdot2^\alpha)\bigr)2^\alpha=c/2>0$.
Apply Theorem~\ref{thm:stretched_ball_weighted_row} with this
value of $a$.
\end{proof}

The constant in~(\ref{eq:stretched_ball_half_rate}) uses the literal
full moment $M_4(c/2,\alpha)\geq1$, by
Theorem~\ref{thm:stretched_moment_one_le}.  It depends only on
$C,B,V,c,\alpha$; there is no factor involving the total number of
labels or sites.  Thus fixed quadratic counting constants give
a bound uniform over finite domains.  Repeated anchor positions
are retained in the label counts.  By
Lemma~\ref{lem:stretched_ball_zero}, all entries at infinite pair
distance vanish before the real-distance convention is used in
the row weights.
